% !TeX root=../main.tex
\chapter{مبانی نظری و تعاریف}
\label{chap:background}

این فصل مفاهیمی را معرفی می‌کند که برای دنبال‌کردن مسئله و روش این پروژه لازم‌اند. متن با فرض آشنایی خواننده با مباحث عمومی دورهٔ کارشناسی مهندسی کامپیوتر نوشته شده است، نه با فرض آشنایی او با روش‌های صوری یا سیستم‌های زمان‌دار. ازاین‌رو، هر موضوع ابتدا به‌صورت شهودی و سپس، هرجا در ادامه به دقت ریاضی نیاز باشد، به‌صورت صوری بیان می‌شود. هدف ارائهٔ یک مرجع جامع برای هریک از این حوزه‌ها نیست؛ خواننده‌ای که با یک موضوع آشناست می‌تواند از بخش مربوط به آن عبور کند.

مسیر فصل از چرایی درستی‌یابی نرم‌افزار آغاز می‌شود و پس از معرفی وارسی مدل و سیستم‌های بلادرنگ، به اتوماتای زمان‌دار و معنای عملیاتی آن می‌رسد. سپس هم‌ارزی‌های رفتاری و تعریف شبیه‌سازی دوسویهٔ ضعیف زمان‌دار ارائه می‌شوند. در پایان نیز ارتباط این تعریف با افراز فضای حالت و کاهش مدل توضیح داده می‌شود.

\section{از آزمون نرم‌افزار تا درستی‌یابی صوری}

نرم‌افزار معمولاً با مجموعه‌ای از آزمون‌های واحد، یکپارچه و سیستمی ارزیابی می‌شود. هر آزمون، اجرای سامانه را برای یک ورودی و شرایط مشخص مشاهده می‌کند و می‌تواند وجود یک خطا را آشکار سازد. بااین‌حال، موفقیت مجموعه‌ای متناهی از آزمون‌ها به‌تنهایی نشان نمی‌دهد که سامانه در همهٔ ورودی‌ها و همهٔ ترتیب‌های اجرای ممکن درست رفتار می‌کند. این محدودیت به‌ویژه در سامانه‌های هم‌روند و واکنشی مهم است؛ زیرا زمان‌بندی متفاوت رویدادها می‌تواند تعداد بسیار زیادی رفتار ممکن پدید آورد \cite{BaierKatoen2008}.

\gls{SoftwareVerification} مفهومی گسترده‌تر از اجرای آزمون‌هاست و بررسی می‌کند که یک مصنوع نرم‌افزاری نیازمندی‌های تعیین‌شده را برآورده می‌کند یا نه. در \gls{FormalVerification}، خود سامانه و ویژگی مورد انتظار با زبان ریاضی توصیف می‌شوند و درستی رابطهٔ میان آن‌ها با استدلال ریاضی یا یک روش الگوریتمی بررسی می‌شود. بنابراین، یک مسئلهٔ درستی‌یابی صوری معمولاً سه جزء دارد: مدلی از سامانه، مشخصه‌ای صوری از رفتار مطلوب یا نامطلوب، و روشی برای اثبات یا رد برقراری آن مشخصه.

نتیجهٔ درستی‌یابی همواره نسبت به مدل و مشخصه‌ای است که تحلیل شده‌اند. اثبات درست‌بودن یک مدل، خودبه‌خود خطای پیاده‌سازی، فرض نادرست محیط، یا نیازمندی‌ای را که در مشخصه نیامده است حذف نمی‌کند. به همین دلیل، آزمون و درستی‌یابی صوری جانشین کامل یکدیگر نیستند: آزمون رفتار پیاده‌سازی واقعی را در اجراهای انتخاب‌شده می‌سنجد، درحالی‌که درستی‌یابی صوری دربارهٔ همهٔ رفتارهای مدل، در محدودهٔ فرض‌های صریح آن، استدلال می‌کند.

دو دستهٔ رایج از ویژگی‌ها عبارت‌اند از ویژگی‌های ایمنی و زنده‌بودن. ویژگی ایمنی بیان می‌کند که «رخداد بدی اتفاق نمی‌افتد»؛ برای نمونه، دو فرایند هم‌زمان وارد یک بخش بحرانی نمی‌شوند. ویژگی زنده‌بودن بیان می‌کند که «رخداد مطلوبی سرانجام اتفاق می‌افتد»؛ برای نمونه، هر درخواست پذیرفته‌شده سرانجام پاسخی دریافت می‌کند. بیان دقیق این ویژگی‌ها ممکن است با منطق‌های زمانی انجام شود، اما جزئیات آن‌ها برای تعریف‌های این فصل لازم نیست \cite{BaierKatoen2008}.

\section{وارسی مدل}

\gls{ModelChecking} یکی از روش‌های خودکار درستی‌یابی صوری برای سامانه‌هایی است که رفتارشان به‌صورت حالت و گذار مدل شده است. اگر $M$ مدل سامانه و $\varphi$ ویژگی مورد نظر باشد، پرسش اصلی به‌شکل
\[
M \models \varphi
\]
نوشته می‌شود: آیا همهٔ رفتارهای مدل $M$ مشخصهٔ $\varphi$ را ارضا می‌کنند؟ الگوریتم وارسی مدل فضای حالت قابل‌دسترسی را به‌طور نظام‌مند بررسی می‌کند. اگر ویژگی برقرار نباشد، ابزار معمولاً یک مثال نقض ارائه می‌کند؛ یعنی اجرایی از مدل که چگونگی نقض ویژگی را نشان می‌دهد \cite{ClarkeGrumbergPeled1999,BaierKatoen2008}.

مزیت وارسی مدل، خودکاربودن و پوشش همهٔ رفتارهای مدلی است که در محدودهٔ تحلیل قرار دارد. مانع اصلی آن \gls{StateSpaceExplosion} است: تعداد حالت‌های مدل ممکن است با افزایش مؤلفه‌ها، متغیرها یا ترتیب‌های ممکن اجرا به‌سرعت رشد کند. در چنین شرایطی، حتی اگر هر گام بررسی ساده باشد، نگهداری و پیمایش همهٔ حالت‌ها پرهزینه یا ناممکن می‌شود. کاهش مدلی که رفتارهای مرتبط با ویژگی مورد بررسی را حفظ کند، یکی از راه‌های مقابله با این مسئله است.

\section{سیستم‌های واکنشی و بلادرنگ}

یک \gls{ReactiveSystem} در تعامل پیوسته با محیط قرار دارد: ورودی دریافت می‌کند، حالت داخلی خود را تغییر می‌دهد و در پاسخ، خروجی یا رویدادی تولید می‌کند. پروتکل‌های ارتباطی، کنترل‌کننده‌ها و سامانه‌های توزیع‌شده نمونه‌هایی از این الگو هستند. در یک سامانهٔ هم‌روند، چند مؤلفه می‌توانند به‌صورت مستقل پیشرفت کنند و ترتیب نسبی کنش‌های آن‌ها بخشی از رفتار سامانه است.

در \gls{RealTimeSystem}، درست‌بودن پاسخ فقط به مقدار منطقی آن وابسته نیست و زمان تولید پاسخ نیز اهمیت دارد. برای نمونه، فرمان درست یک کنترل‌کننده اگر پس از مهلت لازم صادر شود می‌تواند بی‌فایده باشد. در سامانهٔ بلادرنگ سخت، نقض یک مهلت زمانی می‌تواند شکست سامانه محسوب شود؛ در سامانهٔ بلادرنگ نرم، دیرکرد معمولاً کیفیت خدمت را کاهش می‌دهد، اما الزاماً به شکست کامل نمی‌انجامد.

برای تحلیل چنین سامانه‌ای باید گذر زمان نیز در مدل قابل مشاهده باشد. یک مدل صرفاً ترتیبی ممکن است نشان دهد که کنش $a$ پیش از $b$ رخ داده است، اما بیان نمی‌کند فاصلهٔ زمانی میان آن دو چقدر بوده یا آیا یک مهلت رعایت شده است. اتوماتای زمان‌دار این کمبود را با افزودن ساعت و قیدهای زمانی به اتوماتا جبران می‌کند \cite{AlurDill1994,LiebEtAl2025}.

\section{از اتوماتا تا اتوماتای زمان‌دار}

خواننده در درس زبان‌های صوری و نظریهٔ اتوماتا در دورهٔ کارشناسی با ایدهٔ اصلی اتوماتا آشنا شده است: مجموعه‌ای از حالت‌ها و گذارهای برچسب‌دار که نحوهٔ تغییر حالت را مشخص می‌کنند. در این فصل تنها همان بخش از این مفهوم بازخوانی می‌شود که برای مدل‌های زمان‌دار لازم است. در ادبیات اتوماتای زمان‌دار معمولاً به‌جای «حالت کنترلی» از واژهٔ \emph{مکان} و به‌جای یال از واژهٔ \emph{سوییچ} استفاده می‌شود.

اتوماتای زمان‌دار، گراف متناهی اتوماتای کلاسیک را با مجموعه‌ای متناهی از ساعت‌های نامنفی گسترش می‌دهد. همهٔ ساعت‌ها با آهنگی یکسان پیش می‌روند، اما هر گذار می‌تواند تعدادی از آن‌ها را مستقل از بقیه به صفر بازنشانی کند. قیدهای روی ساعت‌ها تعیین می‌کنند یک گذار چه زمانی مجاز است و ماندن در یک مکان تا چه زمانی امکان دارد \cite{AlurDill1994}.

\begin{definition}[اتوماتای زمان‌دار]
\label{def:timed-automaton}
یک اتوماتای زمان‌دار چندتایی
\[
A=(L,l_0,\Sigma,C,I,E)
\]
است که در آن $L$ مجموعه‌ای متناهی از مکان‌ها، $l_0\in L$ مکان آغازین، $\Sigma$ الفبای متناهی کنش‌ها و $C$ مجموعه‌ای متناهی از ساعت‌هاست. تابع $I:L\rightarrow\mathcal{B}(C)$ به هر مکان یک ناوردا نسبت می‌دهد. همچنین
\[
E\subseteq L\times\mathcal{B}(C)\times\Sigma\times 2^C\times L
\]
مجموعهٔ سوییچ‌هاست و $\mathcal{B}(C)$ مجموعهٔ قیدهای ساعتی روی $C$ را نشان می‌دهد. هر عضو $E$ یک سوییچ به‌شکل زیر است:
\[
(l,g,a,R,l').
\]
این سوییچ تنها هنگامی قابل اجراست که شرط نگهبان $g$ برقرار باشد. اجرای آن کنش $a$ را انجام می‌دهد، ساعت‌های مجموعهٔ $R$ را صفر می‌کند و مکان فعال را به $l'$ تغییر می‌دهد. ناوردای $I(l)$ نیز قیدی است که تا هنگام حضور در مکان $l$ باید برقرار بماند \cite{AlurDill1994,LiebEtAl2025}.
\end{definition}

دامنهٔ زمان می‌تواند گسسته، مانند $\mathbb{N}_{\geq 0}$، یا چگال، مانند $\mathbb{R}_{\geq 0}$، باشد. انتخاب دامنه بخشی از مدل‌سازی است و باید با سامانه و روش تحلیل سازگار باشد. در مدل چگال، میان هر دو لحظهٔ متمایز لحظه‌های دیگری وجود دارند؛ در مدل گسسته، زمان با گام‌های شمارا پیش می‌رود.

شکل \ref{fig:timed-automata-examples} شش نمونه از اتوماتای زمان‌دار را نشان می‌دهد. دایره‌ها مکان‌ها هستند. روی یک سوییچ ممکن است کنش، شرط نگهبان و بازنشانی ساعت نوشته شود؛ برای مثال، $x:=0$ ساعت $x$ را هنگام عبور صفر می‌کند. قید نوشته‌شده داخل یک مکان، ناوردای آن مکان است. این نمونه‌ها همچنین نشان می‌دهند که دو اتوماتا با ساختار ظاهری مشابه می‌توانند به‌دلیل تفاوت در نگهبان یا بازنشانی ساعت، رفتار زمانی متفاوتی داشته باشند.

\begin{figure}[htbp]
\centering
\includegraphics[width=\textwidth]{timed-automata-examples.png}
\caption{شش نمونه از اتوماتای زمان‌دار؛ برگرفته از شکل ۱ مقالهٔ لیب و همکاران \cite{LiebEtAl2025}.}
\label{fig:timed-automata-examples}
\end{figure}

\begin{definition}[مقداردهی ساعت]
\label{def:clock-valuation}
برای مجموعهٔ ساعت‌های $C$ و دامنهٔ زمان $\mathbb{T}$، یک \emph{مقداردهی ساعت} تابع $u:C\rightarrow\mathbb{T}$ است. مقدار $u(c)$ زمان سپری‌شده از آخرین بازنشانی ساعت $c$ را نشان می‌دهد. نماد $u+d$ مقداردهی‌ای است که برای هر $c\in C$ مقدار $u(c)+d$ را برمی‌گرداند. اگر $R\subseteq C$ باشد، $[R\mapsto 0]u$ ساعت‌های عضو $R$ را صفر و مقدار بقیهٔ ساعت‌ها را بدون تغییر نگه می‌دارد. همچنین $u\models\phi$ یعنی مقداردهی $u$ قید ساعتی $\phi$ را ارضا می‌کند \cite{LiebEtAl2025}.
\end{definition}

\section{سامانه‌های انتقال برچسب‌دار و رفتار اجرایی}

برای تحلیل یک مدل باید میان نمایش نحوی آن و رفتار حاصل از اجرای آن تمایز گذاشت. اتوماتای زمان‌دار یک نمایش متناهی و فشرده است، اما معنای اجرایی آن مشخص می‌کند از هر وضعیت دقیقاً چه کنش یا چه مقدار تأخیر ممکن است. این معنا با \gls{LabeledTransitionSystem} بیان می‌شود.

\begin{definition}[سامانهٔ انتقال برچسب‌دار]
\label{def:lts}
یک سامانهٔ انتقال برچسب‌دار چندتایی $S=(Q,q_0,\Lambda,\rightarrow)$ است که $Q$ مجموعهٔ حالت‌ها، $q_0\in Q$ حالت آغازین، $\Lambda$ مجموعهٔ برچسب‌ها و $\rightarrow\;\subseteq Q\times\Lambda\times Q$ رابطهٔ گذار است. نماد $q\xrightarrow{\mu}q'$ بیان می‌کند که سامانه می‌تواند با گذار دارای برچسب $\mu$ از $q$ به $q'$ برود.
\end{definition}

یک \emph{اجرا} دنباله‌ای متناوب از حالت‌ها و گذارهاست که از حالت آغازین یا یک حالت مشخص شروع می‌شود. با کنارگذاشتن خود حالت‌ها، دنبالهٔ برچسب‌های این اجرا به دست می‌آید. اگر تنها کنش‌های قابل مشاهده نگه داشته شوند، این دنباله \gls{Trace} نامیده می‌شود. تفاوت میان اجرا و رد مهم است: چند اجرای متفاوت ممکن است رد یکسانی تولید کنند، درحالی‌که انتخاب‌ها و انشعاب‌های داخلی آن‌ها یکسان نباشد.

\begin{definition}[سامانهٔ انتقال برچسب‌دار زمان‌دار]
\label{def:tlts}
فرض کنید $A=(L,l_0,\Sigma,C,I,E)$ یک اتوماتای زمان‌دار روی دامنهٔ زمان $\mathbb{T}$ باشد. \gls{TimedLabeledTransitionSystem} متناظر با آن چندتایی
\[
S_A=(Q,q_0,\Sigma\cup\mathbb{T},\rightarrow)
\]
است که $Q=L\times(C\rightarrow\mathbb{T})$ و $q_0=(l_0,[C\mapsto0])$ است. دو نوع گذار دارد:
\begin{itemize}
\item گذار تأخیر $(l,u)\xrightarrow{d}(l,u+d)$ برای $d\in\mathbb{T}$، مشروط به اینکه ناوردای مکان در طول این تأخیر نقض نشود؛
\item گذار کنش $(l,u)\xrightarrow{a}(l',[R\mapsto0]u)$، هرگاه سوییچ $(l,g,a,R,l')$ وجود داشته باشد، $u\models g$ و مقداردهی مقصد ناوردای $l'$ را ارضا کند.
\end{itemize}
\cite{LiebEtAl2025}
\end{definition}

در نتیجه، یک حالت معنایی فقط نام مکان نیست و مقدار همهٔ ساعت‌ها را نیز در بر می‌گیرد. حتی اتوماتایی با تعداد متناهی مکان می‌تواند سامانهٔ انتقال زمان‌داری با بی‌نهایت حالت داشته باشد. شکل \ref{fig:tlts-examples} بخشی از سامانه‌های انتقال متناظر با سه اتوماتای شکل \ref{fig:timed-automata-examples} را در زمان گسسته نشان می‌دهد. یال‌های با برچسب $1$ گذر یک واحد زمان و یال‌های با برچسب‌های $a$، $b$ و $c$ اجرای کنش را نشان می‌دهند.

\begin{figure}[htbp]
\centering
\includegraphics[width=\textwidth]{tlts-examples.png}
\caption{بخش‌هایی از سامانه‌های انتقال برچسب‌دار زمان‌دار متناظر با $A_1$، $A_3$ و $A_4$؛ برگرفته از شکل ۲ مقالهٔ لیب و همکاران \cite{LiebEtAl2025}.}
\label{fig:tlts-examples}
\end{figure}

\section{هم‌ارزی‌های رفتاری}

برای مقایسهٔ یک مشخصه و پیاده‌سازی، یا برای تشخیص حالت‌های قابل ادغام در یک مدل، باید روشن باشد «رفتار یکسان» دقیقاً به چه معناست. برابری ردها تنها مجموعهٔ دنباله‌های قابل مشاهده را مقایسه می‌کند و ممکن است محل انتخاب میان رفتارهای آینده را نادیده بگیرد. شبیه‌سازی رابطه‌ای جهت‌دار است: هر گام یک سامانه باید در سامانهٔ دیگر قابل تقلید باشد. \gls{Bisimulation} این شرط را در هر دو جهت برقرار می‌کند و در نتیجه ساختار انشعاب رفتار را نیز حفظ می‌کند.

\begin{definition}[شبیه‌سازی دوسویهٔ زمان‌دار قوی]
\label{def:strong-timed-bisimulation}
فرض کنید $S_1=(Q_1,q_{01},\Sigma\cup\mathbb{T},\rightarrow_1)$ و $S_2=(Q_2,q_{02},\Sigma\cup\mathbb{T},\rightarrow_2)$ دو سامانهٔ انتقال زمان‌دار باشند. رابطهٔ $R\subseteq Q_1\times Q_2$ یک شبیه‌سازی دوسویهٔ زمان‌دار قوی است، هرگاه برای هر $(p,q)\in R$ و هر $\mu\in\Sigma\cup\mathbb{T}$ شروط زیر برقرار باشند:
\begin{enumerate}
\item اگر $p\xrightarrow{\mu}_1p'$، حالتی مانند $q'$ وجود داشته باشد که $q\xrightarrow{\mu}_2q'$ و $(p',q')\in R$؛
\item اگر $q\xrightarrow{\mu}_2q'$، حالتی مانند $p'$ وجود داشته باشد که $p\xrightarrow{\mu}_1p'$ و $(p',q')\in R$.
\end{enumerate}
اگر افزون بر این، $(q_{01},q_{02})\in R$ باشد، دو سامانه شبیه‌سازی دوسویهٔ زمان‌دار قوی‌اند \cite{Cerans1993,LiebEtAl2025}.
\end{definition}

واژهٔ «قوی» در این تعریف یعنی هر گذار باید جداگانه و با همان برچسب پاسخ داده شود. اگر مدل گذار داخلی داشته باشد، آن گذار نیز از این قاعده مستثنا نیست. این حساسیت برای مقایسهٔ دقیق دو مدل مفید است، اما هنگام کاهش مدل می‌تواند جزئیات داخلی‌ای را حفظ کند که از دید ناظر اهمیتی ندارند.

\section{کنش داخلی و شبیه‌سازی دوسویهٔ ضعیف زمان‌دار}

همهٔ گذارهای یک مدل الزاماً برای محیط قابل مشاهده نیستند. در این پایان‌نامه، $\varepsilon$ برچسب یک \gls{InternalAction} و $\Sigma$ مجموعهٔ کنش‌های قابل مشاهده است. رویکرد ضعیف اجازه می‌دهد تعدادی گذار داخلی پیرامون یک کنش قابل مشاهده یا در خلال سپری‌شدن زمان نادیده گرفته شوند. بااین‌حال، مقدار کل زمان سپری‌شده باید حفظ شود.

برای بیان تعریف، یک اجرای متناهی را به‌صورت
\[
q_0\xrightarrow{\mu_1}q_1\xrightarrow{\mu_2}\cdots
\xrightarrow{\mu_n}q_n
\]
در نظر بگیرید که هر $\mu_i$ یا یک کنش از $\Sigma_\varepsilon=\Sigma\cup\{\varepsilon\}$ یا مدتی از $\mathbb{R}_{\geq0}$ است. $\operatorname{Untimed}(w)$ دنباله‌ای است که با حذف تأخیرها و کنش‌های $\varepsilon$ از برچسب‌های اجرای $w$ به دست می‌آید؛ $\operatorname{Duration}(w)$ نیز مجموع برچسب‌های زمانی آن اجراست.

\begin{definition}[رابطهٔ گذار ضعیف]
\label{def:weak-transition}
فرض کنید $S=(Q,q_0,\Sigma_\varepsilon,\rightarrow)$ یک سامانهٔ انتقال زمان‌دار باشد که در آن $\rightarrow\;\subseteq Q\times(\Sigma_\varepsilon\cup\mathbb{R}_{\geq0})\times Q$. رابطهٔ گذار ضعیف را با $\overset{\mu}{\Rightarrow}_{W}$ نشان می‌دهیم:
\begin{itemize}
\item برای $d\in\mathbb{R}_{\geq0}$، رابطهٔ $q\overset{d}{\Rightarrow}_{W}q'$ برقرار است، اگر اجرایی $w$ از $q$ به $q'$ وجود داشته باشد و
\[
\operatorname{Untimed}(w)=\varepsilon,
\qquad \operatorname{Duration}(w)=d;
\]
\item برای $a\in\Sigma$، رابطهٔ $q\overset{a}{\Rightarrow}_{W}q'$ برقرار است، اگر اجرایی $w$ از $q$ به $q'$ وجود داشته باشد و
\[
\operatorname{Untimed}(w)=a,
\qquad \operatorname{Duration}(w)=0.
\]
\end{itemize}
بنابراین، یک گام ضعیفِ کنش می‌تواند گذارهای داخلی پیش و پس از دقیقاً یک کنش قابل مشاهده را در خود جمع کند، بدون آنکه زمانی بگذرد. یک گام ضعیفِ زمانی نیز می‌تواند چند گذار زمانی و داخلی را ترکیب کند، به‌شرط آنکه هیچ کنش قابل مشاهده‌ای رخ ندهد و مجموع تأخیرها دقیقاً $d$ باشد \cite{ReynierSangnier2009}.
\end{definition}

\begin{definition}[شبیه‌سازی دوسویهٔ ضعیف زمان‌دار]
\label{def:weak-timed-bisimulation}
فرض کنید
$S_1=(Q_1,q_{01},\Sigma_\varepsilon,\rightarrow_1)$ و
$S_2=(Q_2,q_{02},\Sigma_\varepsilon,\rightarrow_2)$
دو سامانهٔ انتقال زمان‌دار باشند. رابطهٔ $R\subseteq Q_1\times Q_2$ یک شبیه‌سازی دوسویهٔ ضعیف زمان‌دار است، هرگاه $(q_{01},q_{02})\in R$ و برای هر $(p,q)\in R$ و هر $\mu\in\Sigma\cup\mathbb{R}_{\geq0}$ شروط زیر برقرار باشند:
\begin{enumerate}
\item اگر $p\overset{\mu}{\Rightarrow}_{W,1}p'$، حالتی مانند $q'$ وجود داشته باشد که $q\overset{\mu}{\Rightarrow}_{W,2}q'$ و $(p',q')\in R$؛
\item اگر $q\overset{\mu}{\Rightarrow}_{W,2}q'$، حالتی مانند $p'$ وجود داشته باشد که $p\overset{\mu}{\Rightarrow}_{W,1}p'$ و $(p',q')\in R$.
\end{enumerate}
وجود چنین رابطه‌ای را با $S_1\approx_W S_2$ نشان می‌دهیم \cite{ReynierSangnier2009}.
\end{definition}

تفاوت اصلی تعریف ضعیف با تعریف قوی در این است که سامانهٔ پاسخ‌دهنده مجبور نیست تک‌تک گام‌های داخلی سامانهٔ نخست را بازتولید کند. آنچه باید یکسان بماند، کنش قابل مشاهده، زمان کل سپری‌شده و امکان ادامهٔ رفتار از حالت‌های مقصد است. این نکته اجازه می‌دهد حالت‌هایی که تنها در جزئیات داخلی تفاوت دارند در یک ردهٔ رفتاری قرار گیرند، بی‌آنکه تمایز زمانی قابل مشاهده از میان برود.

\section{بازنمایی متناهی زمان}

در زمان چگال، حتی یک اتوماتای زمان‌دار کوچک می‌تواند بی‌نهایت مقداردهی ساعت و در نتیجه بی‌نهایت حالت معنایی داشته باشد. برای تحلیل الگوریتمی باید مجموعه‌های نامتناهی از مقداردهی‌ها به‌شکل متناهی نمایش داده شوند. دو بازنمایی شناخته‌شده برای این منظور گراف ناحیه و گراف \lr{zone} هستند \cite{AlurDill1994,HenzingerEtAl1994,LiebEtAl2025}.

در گراف ناحیه، مقداردهی‌هایی که همهٔ قیدهای مرتبط را یکسان ارزیابی می‌کنند و ترتیب عبور بخش کسری ساعت‌هایشان از مرزهای صحیح مشابه است در یک کلاس هم‌ارزی قرار می‌گیرند. تعداد این کلاس‌ها برای یک اتوماتای معین متناهی است؛ ازاین‌رو گراف ناحیه امکان تصمیم‌گیری دربارهٔ برخی مسائل را فراهم می‌کند، اما با بزرگ‌شدن ثابت‌های زمانی یا تعداد ساعت‌ها می‌تواند بسیار بزرگ شود \cite{AlurDill1994,Cerans1993}.

یک \lr{zone} مجموعه‌ای از مقداردهی‌های ساعت است که با یک هم‌بندی از قیدهای ساعتی نمایش داده می‌شود. گراف‌های \lr{zone} معمولاً چند ناحیه را در یک حالت نمادین جمع می‌کنند و در عمل از گراف ناحیه فشرده‌ترند. ماتریس کران تفاضل یا \lr{DBM} نمایش متداول این قیدهاست \cite{HenzingerEtAl1994,BengtssonYi2004}. بااین‌حال، فشرده‌سازی ممکن است اطلاعاتی مانند اثر یک بازنشانی ساعت بر انشعاب‌های آینده را پنهان کند؛ بنابراین مناسب‌بودن یک بازنمایی نمادین به ویژگی رفتاری مورد بررسی وابسته است. لیب و همکاران برای وارسی شبیه‌سازی دوسویهٔ زمان‌دار، گراف‌های \lr{zone} را با ساعت‌های مجازی گسترش داده‌اند تا بخشی از این اطلاعات حفظ شود \cite{LiebEtAl2025}.

\section{هم‌ارزی، افراز و کاهش مدل}

شبیه‌سازی دوسویه یک رابطهٔ هم‌ارزی است؛ یعنی حالت‌ها را می‌توان بر اساس آن به رده‌های جدا از هم تقسیم کرد. چنین تقسیم‌بندی‌ای \gls{Partition} نام دارد. در هر بلوک افراز، حالت‌هایی قرار می‌گیرند که مطابق تعریف انتخاب‌شده از دید رفتاری تمایزپذیر نیستند.

با جایگزین‌کردن هر بلوک با یک حالت و انتقال گذارها میان بلوک‌ها، مدل خارج‌قسمتی به دست می‌آید. اگر افراز نسبت به شبیه‌سازی دوسویهٔ ضعیف زمان‌دار پایدار باشد، ادغام حالت‌های یک بلوک کنش‌های قابل مشاهده، زمان کل و ساختار انشعاب مورد نظر تعریف را حفظ می‌کند. هدف کاهش مدل، ساختن چنین نمایش کوچک‌تری است؛ نه حذف دلخواه حالت‌ها. الگوریتم‌های پالایش افراز معمولاً از یک افراز اولیه شروع می‌کنند و هر بلوکی را که اعضایش نتوانند گام‌های ضعیف یکدیگر را پاسخ دهند می‌شکنند. این فرایند تا رسیدن به افرازی پایدار ادامه می‌یابد.

این زنجیرهٔ مفهومی جایگاه مسئله را روشن می‌کند: وارسی مدل با رشد فضای حالت محدود می‌شود؛ مدل‌های بلادرنگ علاوه بر کنش، گذر زمان را نیز ثبت می‌کنند؛ و شبیه‌سازی دوسویهٔ ضعیف زمان‌دار معیاری فراهم می‌کند تا جزئیات داخلی پنهان شوند، درحالی‌که رفتار زمانی قابل مشاهده حفظ می‌شود. فصل‌های بعد باید مشخص کنند مدل ورودی این پروژه دقیقاً چه ساختاری دارد، تعریف بالا چگونه در الگوریتم به کار رفته است و کاهش حاصل با چه شواهدی ارزیابی می‌شود.
